\section*{Введение}
\addcontentsline{toc}{section}{Введение}

В отчёте представлено описание лабораторных работ №1–№5 по дисциплине «Теория графов». 
В ходе выполнения работы необходимо реализовать программу, которая строит связный ациклический ориентированный граф в соответствии с заданным распределением (распределение Чампернауна), а также реализовать требуемые методы и алгоритмы над графом.

\subsection*{Лабораторная работа №1}
\begin{enumerate}
	\item Сформировать случайным образом связный ациклический ориентированный граф в соответствии с распределением Чампернауна (параметры распределения задаются как константы) с вводимым пользователем количеством вершин. Генерация выполняется по степеням вершин. В зависимости от дальнейших заданий используется ориентированный или неориентированный граф, полученный из ориентированного.
	\item Вычислить эксцентриситеты вершин, определить центр графа и диаметральные вершины.
	\item Реализовать метод Шимбелла для весовой матрицы, сгенерированной в соответствии с заданным распределением. В ответе представить минимальный и/или максимальный путь в виде матрицы. Весовая матрица генерируется с возможностью выбора пользователем: только положительные веса, только отрицательные или смешанные.
	\item Определить возможность построения маршрута от одной заданной точки до другой (вершины вводит пользователь) и указать количество таких маршрутов.
	\item Реализовать меню для выбора действий.
\end{enumerate}

\subsection*{Лабораторная работа №2}
\begin{enumerate}
	\item Для графов, сгенерированных в первой работе, реализовать алгоритм поиска точек сочленения.
	\item Сгенерировать весовую матрицу (положительную или с отрицательными весами – на выбор пользователя) и найти кратчайший путь для двух выбранных пользователем вершин с помощью классического алгоритма Дейкстры. Результат содержит расстояние, сам путь в виде последовательности вершин и вектор расстояний.
	\item Сравнить скорость работы реализованных алгоритмов по количеству итераций.
\end{enumerate}

\subsection*{Лабораторная работа №3}
\begin{enumerate}
	\item На основе сгенерированного ориентированного графа случайным образом получить матрицы пропускных способностей и стоимости.
	\item Для полученного графа найти максимальный поток по алгоритму Форда–Фалкерсона.
	\item Вычислить поток минимальной стоимости. Величина потока задаётся как $\lceil \frac{2}{3} \cdot \max \rceil$, где $\max$ – значение максимального потока. Для этого использовать ранее реализованный алгоритм Дейкстры.
\end{enumerate}

\subsection*{Лабораторная работа №4}
\begin{enumerate}
	\item Для неориентированных графов, сгенерированных в первой работе, найти число остовных деревьев с помощью матричной теоремы Кирхгофа.
	\item Построить минимальный по весу остов для сгенерированного неориентированного взвешенного графа, используя алгоритм Борувки. Полученный остов закодировать с помощью кода Прюфера и декодировать его, обязательно сохраняя веса.
	\item Реализовать на исходном графе и полученном остове (по выбору пользователя) алгоритм нахождения минимальной раскраски графа.
\end{enumerate}

\subsection*{Лабораторная работа №5}
\begin{enumerate}
	\item Для неориентированных графов, сгенерированных в первой работе, проверить, является ли граф эйлеровым. Если нет – модифицировать граф (логировать изменения) и построить эйлеров цикл.
	\item Используя кратчайший остов, построенный в лабораторной работе №4, для получить фундаментальную систему разрезов (вывести на экран). На основе фундаментальных разрезов с помощью операции симметрической разности получить остальные разрезы (выбранные разрезы вводит пользователь).
\end{enumerate}